\chapter{Architecture mat\'erielle}

Ce chapitre d\'ecrit en d\'etail les composants mat\'eriels du syst\`eme Andon, depuis les n\oe{}uds de d\'etection jusqu'\`a l'infrastructure r\'eseau.

\section{Le n\oe{}ud Andon ESP32}

\subsection{Choix du microcontr\^oleur}

L'ESP32 (Espressif Systems) a \'et\'e s\'electionn\'e comme microcontr\^oleur principal pour chaque n\oe{}ud Andon. Ce choix est motiv\'e par :

\begin{itemize}
    \item Connectivit\'e Wi-Fi 802.11 b/g/n et Bluetooth 4.2/BLE int\'egr\'ees.
    \item Processeur dual-core Xtensa LX6 \`a 240 MHz avec 520 KB de RAM.
    \item 34 broches GPIO dont 16 ADC (convertisseur analogique-num\'erique).
    \item Interface I$^2$C et SPI pour l'ajout de capteurs externes.
    \item Consommation r\'eduite : 80 mA en mode Wi-Fi actif.
    \item Prix unitaire comp\'etitif ($\approx$ 4EUR).
    \item Soutien Arduino et MicroPython, facilitant le d\'eveloppement.
\end{itemize}

\subsection{Composants du n\oe{}ud}

Chaque n\oe{}ud Andon est compos\'e des \'el\'ements suivants :

\begin{table}[htbp]
\centering
\renewcommand{\arraystretch}{1.3}
\begin{tabular}{l l l}
\toprule
\textbf{Composant} & \textbf{R\'ef\'erence} & \textbf{Fonction} \\
\midrule
Microcontr\^oleur & ESP32-DevKitC & Processeur, Wi-Fi, BT \\
Bouton poussoir & KY-004 & Signalisation Andon \\
LED verte & LED 5mm 3V & Indicateur ligne OK \\
LED rouge & LED 5mm 3V & Indicateur panne \\
R\'esistances & 220 $\Omega$ (x2) & Limitation courant LED \\
R\'esistance pull-down & 10 k$\Omega$ & Stabilisation signal bouton \\
Alimentation & USB 5V / adaptateur & Alimentation du n\oe{}ud \\
Bo\^itier & Bo\^itier plastique IP54 & Protection du circuit \\
\bottomrule
\end{tabular}
\caption{Composants du n\oe{}ud Andon ESP32}
\label{tab:composants_esp32}
\end{table}

\subsection{Sch\'ema de câblage}

Le sch\'ema de câblage d\'efinit les connexions \'electriques du n\oe{}ud Andon :

\begin{itemize}
    \item Le bouton poussoir est connect\'e entre la broche \textbf{GPIO4} et la masse (GND), avec une r\'esistance de pull-down de 10 k$\Omega$ pour garantir un \'etat logique bas d\'efini au repos.
    \item La \textbf{LED verte} est connect\'e entre la broche \textbf{GPIO16} et la masse, s\'erie d'une r\'esistance de 220 $\Omega$.
    \item La \textbf{LED rouge} est connect\'e entre la broche \textbf{GPIO17} et la masse, s\'erie d'une r\'esistance de 220 $\Omega$.
    \item L'alimentation est fournie via le connecteur USB du module ESP32 (5V).
\end{itemize}

\figplaceholder{Sch\'ema de câblage du n\oe{}ud Andon}{fig:schema_cablage_esp32}

\subsection{Fonctionnement du circuit}

Le circuit fonctionne selon le principe suivant :

\begin{enumerate}
    \item Au repos, la broche GPIO4 est \`a l'\`etat bas (0V) gr\^ace \`a la r\'esistance pull-down.
    \item Lorsque l'op\'erateur appuie sur le bouton, le signal passe \`a l'\`etat haut (3,3V).
    \item L'ESP32 d\'etecte le front montant par interruption logicielle.
    \item Apr\`es d\'ebounce (200 ms), l'ESP32 g\'en\`ere un message JSON.
    \item Le message est publi\'e sur le topic MQTT correspondant.
    \item L'ESP32 allume la LED rouge en attendant la confirmation du serveur.
\end{enumerate}

\section{Infrastructure r\'eseau}

\subsection{R\'eseau Wi-Fi}

Le r\'eseau Wi-Fi de l'atelier assure la communication entre les n\oe{}uds ESP32 et le broker MQTT :

\begin{table}[htbp]
\centering
\renewcommand{\arraystretch}{1.3}
\begin{tabular}{l l}
\toprule
\textbf{Param\`etre} & \textbf{Valeur} \\
\midrule
Standard & Wi-Fi 802.11 b/g/n (2,4 GHz) \\
Protocole de s\'ecurit\'e & WPA2-Enterprise \\
Port d\'edi\'e aux IoT & VLAN s\'epar\'e \\
Filtrage & Par adresse MAC \\
Bande passante & 1 Mb/s d\'edi\' par n\oe{}ud \\
Port\'ee & 30--50 m (atelier) \\
\bottomrule
\end{tabular}
\caption{Param\`etres du r\'eseau Wi-Fi}
\label{tab:wifi_params}
\end{table}

\subsection{R\'eseau filaire}

Le serveur backend est connect\'e au r\'eseau filaire de l'entreprise :

\begin{itemize}
    \item Connectivit\'e Ethernet Gigabit.
    \item Adresse IP statique dans le sous-r\'eseau d\'edi\'e.
    \item Pare-feu configur\'e pour limiter les acc\`es aux ports n\'ecessaires (1883 MQTT, 3000 API, 5432 PostgreSQL).
    \item Pas de connexion Internet directe pour des raisons de s\'ecurit\'e.
\end{itemize}

\section{Alimentation et fiabilit\'e}

\subsection{Alimentation des n\oe{}uds}

Chaque n\oe{}ud Andon est aliment\'e par :

\begin{itemize}
    \item Un adaptateur secteur 5V/2A branch\'e sur une prise \`a proximit\'e du poste.
    \item En secours : un module de backup batteries (autonomie 4h).
\end{itemize}

\subsection{Robustesse industrielle}

Les n\oe{}uds doivent r\'esister aux conditions de l'atelier :

\begin{itemize}
    \item Temp\'erature ambiante : 15$^\circ$C \`a 35$^\circ$C.
    \item Humidit\'e relative : 30\% \`a 80\%.
    \item Poussi\`ere : bo\^itier IP54 recommand\'e.
    \item Vibrations : fixation solide du bo\^itier au poste de travail.
    \item Perturbations \'electromagn\'etiques : blindage du bo\^itier m\'etallique.
\end{itemize}

\section{Dimensionnement du syst\`eme}

\subsection{Nombre de n\oe{}uds}

Le syst\`eme est dimensionn\'e pour l'atelier Test \'Electrique de \sews :

\begin{table}[htbp]
\centering
\renewcommand{\arraystretch}{1.3}
\begin{tabular}{l c}
\toprule
\textbf{Param\`etre} & \textbf{Valeur} \\
\midrule
Nombre de lignes & 8 \\
N\oe{}uds par ligne & 1 \\
N\oe{}uds total & 8 \\
Boutons par ligne & 1 \\
LEDs par ligne & 2 (rouge + verte) \\
Co\^ut par n\oe{}ud & $\approx$ 12EUR \\
Co\^ut total capteurs & $\approx$ 96EUR \\
\bottomrule
\end{tabular}
\caption{Dimensionnement des n\oe{}uds}
\label{tab:dimensionnement}
\end{table}

\section{Synth\`ese du chapitre}

L'architecture mat\'erielle du syst\`eme repose sur des n\oe{}uds ESP32 autonomes, chacun \'e\'equip\'e d'un bouton Andon et de deux LED indicatrices. Le r\'eseau Wi-Fi de l'atelier assure la transmission des donn\'ees vers un serveur backend. Le co\^ut total du mat\'eriel est tr\`es comp\'etitif ($\approx$ 96EUR pour 8 n\oe{}uds), d\'emontrant la pertinence de cette approche par rapport aux solutions commerciales.
